\chapter{Mixed Set B}

Twenty problems, no order, no announced methods. Target time: 75 minutes.

\section{Worked Example}

Again the example is not one of the twenty. This one is a constrained
optimisation, one of the most frequently examined topics, worked
with every check that the marker expects to see.

\begin{worked}{Maximise $f(x,y)=\ln x+2\ln y$ subject to $x+y=12$, with $x,y>0$: solve, certify, and interpret the multiplier}
\wstep{Step 1. Recognise the type before writing anything.}
Two variables, one equality constraint, a differentiable objective: this is
Lagrange. The constraint is linear and could be substituted instead --- keep
that in reserve as the check.

\wstep{Step 2. Set up and differentiate.}
\[
  \mathcal L=\ln x+2\ln y+\lambda(12-x-y),
\]
\[
  \mathcal L_{x}=\frac1x-\lambda=0,
  \qquad
  \mathcal L_{y}=\frac2y-\lambda=0,
  \qquad
  x+y=12 .
\]

\wstep{Step 3. Eliminate $\lambda$.}
\[
  \frac1x=\frac2y
  \quad\Longrightarrow\quad
  y=2x .
\]
The optimal bundle puts twice as much into $y$ as into $x$ --- the ratio of
the coefficients $2$ and $1$ in the objective, exactly as Cobb--Douglas
shares behave, and $\ln$-linear objectives are Cobb--Douglas in disguise.

\wstep{Step 4. Impose the constraint.}
\[
  x+2x=12
  \quad\Longrightarrow\quad
  \boxed{x^{*}=4,\quad y^{*}=8,\quad \lambda^{*}=\tfrac14 }
\]
\[
  f^{*}=\ln4+2\ln8=\ln4+\ln64=\ln256\approx5.545 .
\]

\wstep{Step 5. Certify it with the bordered Hessian.}
\[
  \mathcal L_{xx}=-\frac{1}{x^{2}}=-\tfrac{1}{16},
  \qquad
  \mathcal L_{yy}=-\frac{2}{y^{2}}=-\tfrac{1}{32},
  \qquad
  \mathcal L_{xy}=0,
  \qquad
  g_{x}=g_{y}=1 .
\]
\[
  \overline H=\begin{pmatrix}0&1&1\\1&-\tfrac{1}{16}&0\\1&0&-\tfrac{1}{32}\end{pmatrix},
  \qquad
  \det\overline H=\tfrac{1}{32}+\tfrac{1}{16}=\tfrac{3}{32}>0 ,
\]
which for $n=2$, $m=1$ means a constrained \textbf{maximum}. \checkmark

\wstep{Step 6. Check again by substitution, which is available here.}
$y=12-x$ gives $h(x)=\ln x+2\ln(12-x)$, so
\[
  h'(x)=\frac1x-\frac{2}{12-x}=0
  \;\Longrightarrow\;
  12-x=2x
  \;\Longrightarrow\;
  x=4 ,
\]
\[
  h''(x)=-\frac{1}{x^{2}}-\frac{2}{(12-x)^{2}}<0\ \text{everywhere:\ concave, so the maximum is global.}
\]
\wtrap{Two independent routes to the same answer is the best use of spare
examination time --- far better than re-reading the first one. And here the
substitution route delivers something the bordered Hessian does not: a
\emph{global} conclusion, because $h''<0$ on the whole feasible interval,
not merely at the candidate point.}

\wstep{Step 7. Verify the multiplier's meaning.}
With a general budget $M$ in place of $12$: $x^{*}=M/3$, $y^{*}=2M/3$, so
\[
  V(M)=\ln\frac M3+2\ln\frac{2M}{3}=3\ln M+\text{constant},
  \qquad
  \dd{V}{M}=\frac3M .
\]
At $M=12$ this is $\tfrac14=\lambda^{*}$. \checkmark\ One more unit of budget
raises the maximised objective by about $0.25$.

\wstep{Step 8. The checks a marker looks for.}
\begin{center}\footnotesize
\begin{tabular}{ll}
  \hline
  check & here\\
  \hline
  constraint satisfied & $4+8=12$ \checkmark\\
  domain respected & $x,y>0$ \checkmark\\
  second-order condition & $\det\overline H>0$ \checkmark\\
  multiplier sign & $\lambda>0$: relaxing helps \checkmark\\
  interpretation stated & $\lambda=\partial V/\partial M$ \checkmark\\
  \hline
\end{tabular}
\end{center}
Every one of these is a line, and together they are usually worth more marks
than the algebra that produced the numbers.
\end{worked}

\section{Problems}

\begin{problems}
  \item Simplify $\dfrac{x^{2}-4}{x^{2}-x-6}$ and state where the
        simplification is valid.
  \item Solve $\log_{2}x+\log_{2}(x-2)=3$.
  \item Give the range of $f(x)=\dfrac{x^{2}+1}{x}$ on $x>0$.
  \item Write a system of inequalities for the part of the elliptical disc
        $\dfrac{x^{2}}{4}+\dfrac{y^{2}}{9}\le1$ lying in the second
        quadrant.
  \item Two functions have concentric circles as level curves. In one the
        circles are equally spaced, in the other they crowd outwards. Which
        is $\sqrt{x^{2}+y^{2}}$ and which is $x^{2}+y^{2}$?
  \item Evaluate $\displaystyle\lim_{x\to0^{+}}x^{\sin x}$.
  \item Evaluate $\displaystyle\lim_{x\to2}\frac{x^{3}-8}{x^{2}-4}$.
  \item For $f(x)=\dfrac{\mathrm{e}^{x}}{x}$: give the domain, the
        asymptotes, the monotonicity and any extremum.
  \item Solve $2x+y=7$, \;$4x+2y=14$, and describe the solution set.
  \item Compute
        $\det\begin{pmatrix}1&2&0\\0&3&1\\2&0&4\end{pmatrix}$.
  \item Solve $\begin{pmatrix}2&1\\1&3\end{pmatrix}
        \mathbf{x}=\binom{5}{10}$.
  \item Give $\grad f(1,2)$ for $f(x,y)=x^{2}y^{3}$.
  \item Give the directional derivative of the same $f$ at $(1,2)$ in the
        direction of $(3,4)^{T}$.
  \item Use the total differential to estimate $f(1.1,1.9)$ for the same
        $f$, and compare with the exact value.
  \item Find $\dd{y}{x}$ at $(1,2)$ on the curve $x^{2}+xy+y^{2}=7$.
  \item Minimise $2x+3y$ subject to $xy=54$, with $x,y>0$.
  \item Maximise $\ln x+y$ subject to $x+y\le4$, $x>0$, $y\ge0$.
  \item Maximise $4x+3y$ subject to $x+y\le10$, $2x+y\le16$, $x,y\ge0$.
  \item What is $1000$ received in five years worth today at $6\%$ annual
        interest?
  \item Sum the series $3+\tfrac32+\tfrac34+\tfrac38+\cdots$.
\end{problems}

\section{Solutions}

\begin{solutions}
  \item \textbf{Factor and cancel.}
        $\dfrac{(x-2)(x+2)}{(x-3)(x+2)}=\dfrac{x-2}{x-3}$, valid for
        $x\neq-2$ and $x\neq3$. The point $x=-2$ remains excluded from the
        domain even though the factor cancelled --- a hole, not a pole.
  \item \textbf{Logarithm law, then check the domain.}
        $\log_{2}\bigl(x(x-2)\bigr)=3\Rightarrow x^{2}-2x=8$, so
        $x^{2}-2x-8=0$ and $x=4$ or $x=-2$. The domain requires $x>2$, so
        $x=-2$ is rejected: \;$x=4$.
  \item \textbf{Range via AM--GM (or calculus).}
        $f(x)=x+\dfrac1x$, and for $x>0$, $x+1/x\ge2$ with equality at
        $x=1$. As $x\to0^{+}$ and $x\to\infty$ the function grows without
        bound, so $W=[2,\infty)$.
  \item \textbf{Region.} $\dfrac{x^{2}}{4}+\dfrac{y^{2}}{9}\le1$,
        \;$x\le0$, \;$y\ge0$.
  \item \textbf{Contour spacing.} Equally spaced circles belong to the
        \textbf{cone} $\sqrt{x^{2}+y^{2}}$, whose level $c$ has radius $c$.
        Circles crowding outwards belong to the \textbf{paraboloid}
        $x^{2}+y^{2}$, whose level $c$ has radius $\sqrt{c}$: the surface
        steepens, so equal steps in height need ever smaller steps in
        radius.
  \item \textbf{Exponential form, take logs.}
        $\ln y=\sin x\ln x$, and since $\sin x\sim x$ near $0$ this behaves
        like $x\ln x\to0$. Hence the limit is $\mathrm{e}^{0}=1$.
  \item \textbf{L'H\^opital (or factor).}
        $\dfrac{3x^{2}}{2x}\to\dfrac{12}{4}=3$. Factoring gives the same:
        $\dfrac{(x-2)(x^{2}+2x+4)}{(x-2)(x+2)}\to\dfrac{12}{4}$.
  \item \textbf{Curve discussion.} Domain $\R\setminus\{0\}$. As
        $x\to0^{+}$, $f\to+\infty$; as $x\to0^{-}$, $f\to-\infty$: a pole
        at $x=0$. As $x\to-\infty$, $f\to0^{-}$, so $y=0$ is a horizontal
        asymptote on the left only; as $x\to+\infty$, $f\to+\infty$. And
        $f'=\dfrac{\mathrm{e}^{x}(x-1)}{x^{2}}$, negative for $x<1$
        (excluding $0$) and positive for $x>1$: decreasing, then
        increasing, with a local (and global, on $x>0$) minimum at
        $(1,\mathrm{e})$.
  \item \textbf{Elimination.} The second equation is exactly twice the
        first, so it carries no information: one pivot, two unknowns.
        \textbf{Infinitely many} solutions, $y=7-2x$ with $x$ free ---
        the line itself.
  \item \textbf{Cofactor expansion along the first row.}
        $1\cdot(3\cdot4-1\cdot0)-2\cdot(0\cdot4-1\cdot2)+0
        =12+4=\mathbf{16}$.
  \item \textbf{Inverse or elimination.} $\det=6-1=5$, so
        $\mathbf{x}=\dfrac15\begin{pmatrix}3&-1\\-1&2\end{pmatrix}
        \binom{5}{10}=\dfrac15\binom{5}{15}=\binom{1}{3}$.
        (Check: $2+3=5$ \checkmark, $1+9=10$ \checkmark.)
  \item \textbf{Partial derivatives.} $f_{x}=2xy^{3}$ and
        $f_{y}=3x^{2}y^{2}$, so at $(1,2)$:
        $\grad f=(2\cdot8,\;3\cdot4)^{T}=(16,12)^{T}$.
  \item \textbf{Directional derivative --- normalise first.}
        $\|(3,4)\|=5$, so $\mathbf{u}=(0.6,0.8)^{T}$ and
        $D_{\mathbf{u}}f=(16,12)\cdot(0.6,0.8)=9.6+9.6=\mathbf{19.2}$.
  \item \textbf{Total differential.} $f(1,2)=1\cdot8=8$ and
        $\mathrm{d}f=16\,\mathrm{d}x+12\,\mathrm{d}y
        =16(0.1)+12(-0.1)=0.4$, so $f(1.1,1.9)\approx8.4$. Exactly
        $1.1^{2}\cdot1.9^{3}=1.21\cdot6.859=8.299$; the linear estimate is
        high by about $0.1$, the second-order error.
  \item \textbf{Implicit differentiation.} $2x+y+xy'+2yy'=0$, so
        $y'=-\dfrac{2x+y}{x+2y}$. At $(1,2)$ --- which does lie on the
        curve, since $1+2+4=7$ --- this is
        $-\dfrac{4}{5}$.
  \item \textbf{Lagrange (cost minimisation).}
        $2=\lambda y$ and $3=\lambda x$ give $\dfrac{x}{y}=\dfrac32$, so
        $x=1.5y$. Then $1.5y^{2}=54$, $y^{2}=36$, $y=6$ and $x=9$. Minimal
        cost $18+18=36$ --- the two expenditures are equal, as always for
        a product constraint with linear costs.
  \item \textbf{Kuhn--Tucker.} With
        $\mathcal{L}=\ln x+y+\lambda(4-x-y)$: $1/x=\lambda$ and
        $1=\lambda$, so $\lambda=1$ and $x=1$. The budget binds, so $y=3$,
        which satisfies $y\ge0$. Optimum $(1,3)$ with value
        $\ln1+3=3$.
  \item \textbf{Linear programming.} The two constraints meet where
        $x+y=10$ and $2x+y=16$, i.e. $(6,4)$. Vertices $(0,0)$, $(8,0)$,
        $(6,4)$, $(0,10)$ give objective values $0$, $32$, $36$, $30$: the
        maximum is $\mathbf{36}$ at $(6,4)$, where both constraints bind.
  \item \textbf{Discounting.}
        $\dfrac{1000}{1.06^{5}}=\dfrac{1000}{1.33823}\approx747.26$.
  \item \textbf{Geometric series.} First term $a=3$, ratio
        $r=\tfrac12$, so the sum is $\dfrac{3}{1-\tfrac12}=6$.
\end{solutions}
